

%% 
%% This is a LaTeX document generated by automatic conversion of a
%% Mathematica notebook using Mathematica 4.0.
%% 
%% This document uses special macros that are defined in a LaTeX 
%% package file with a name of the form notebook*.sty.  Appropriate 
%% commands for loading the package have been included in this document.
%% The LaTeX package files can be found in a directory with path:
%% 
%% /usr/local/mathematica/SystemFiles/IncludeFiles/TeX/
%% 
%% The LaTeX package used in this document may require that your 
%% LaTeX implementation support fonts other than the standard Computer 
%% Modern fonts that are used by LaTeX.  See the Additional Information 
%% section under the TeXSave entry in the Mathematica Reference Guide 
%% for more information.
%% 
%% Implementation-specific instructions for installing TeX-related files
%% can be found at this URL.
%% 
%% http://www.wolfram.com/support/FrontEnds/Export/TeX
%% 
%% 

\documentclass{article}
\usepackage{notebook2e, latexsym}


\begin{document}

\Title{Se\NTilde{}ales }
\Subtitle{Segunda clase}
\Subsubtitle{EE013 Se\NTilde{}ales y Sistemas}
\Subsubtitle{Daniel W. Berns

Universidad Nacional de la Patagonia San Juan Bosco

}
Esto carga algunas funciones necesarias.

\dispSFinmath{<<\Mvariable{SignalPlot`}}
\dispSFinmath{<<\Mvariable{Graphics`Graphics`}}
\begin{CellGroup}
\Section*{1 Se\NTilde{}ales continuas}


A small number of basic signals are of fundamental importance in the study  of linear systems. This list includes the families of sinusoidal and
 exponential signals, and such special signals as the unit impulse, unit step, and the unit pulse. In this section, the sinusoidal and exponential
signals will  be considered, while special signals will be covered later. Continuous-time  signals are typically defined over an infinite (-\(\infty
\), \(+\)\(\infty \))  or semi-infinite [0,\(+\)\(\infty \)) domain of the real line. In general, the signal is simply a function of one or more
real independent variable(s).

\begin{CellGroup}
\Subsection*{Real-valued sinusoidal signals}
One of the most important signals in linear signal and system theory is  the sinusoid. Two continuous-time sinusoids with different periods are displayed
below.


\end{CellGroup}
\begin{CellGroup}
\dispSFinmath{\MathBegin{MathArray}{l}
\Muserfunction{DisplayTogetherArray}[\{\Mfunction{Plot}[\sin [2\multsp \pi \multsp t],\{t,0,2\},  \\
\noalign{\vspace{0.5ex}}
\hspace{3.em} \Mvariable{PlotLabel}->\multsp ,\Mvariable{AxesLabel}\rightarrow \{t,\Mvariable{Seno(2\multsp Pi\multsp t)}\}],\Mfunction{Plot}[\sin
[4\multsp \pi \multsp t],  \\
\noalign{\vspace{0.5ex}}
\hspace{3.em} \{t,0,2\},\Mvariable{PlotLabel}->\multsp ,\Mvariable{AxesLabel}\rightarrow \{t,\Mvariable{Seno(4\multsp Pi\multsp t)}\}]\}];\\
\MathEnd{MathArray}}
\mathGraphic{lecture02.tex_gr1.eps}

\end{CellGroup}
 A period of a signal \inlineTFinmath{f(t)} is a value {\itshape T,} such that \inlineTFinmath{f(t)\multsp =\multsp f(t+T)}. The period of a sinusoidal
signal of the form \inlineTFinmath{\sin({{\omega }_0}\multsp t)}, where \inlineTFinmath{{{\omega }_0}} is a real value, is given by



\inlineTFinmath{\sin({{\omega }_0}\multsp (t+T))\multsp =\multsp \sin({{\omega }_0}t\multsp +\multsp {{\omega }_0}\multsp T)\multsp =\multsp \sin({{\omega
}_0}t\multsp +\multsp 2\pi )}

By comparing the rightmost two terms in (1), we get the following relationship for the period \inlineTFinmath{T}:



\inlineTFinmath{T=\multsp \frac{2\pi }{{{\omega }_0}}}.

The inverse of the period is called the {\itshape frequency} (strictly speaking, it is called the {\itshape fundamental frequency}) of the sinusoid.
Observe that for a sinusoid, the frequency defines the number of periods of the signal in a unit of time. It is also directly related to the number
of {\itshape zero-crossings} of the sinusoid in a unit of time. The units of frequency are either hertz (Hz) or radians per second depending on the
exact relationship to the period. Units of Hz are used for frequency defined as \inlineTFinmath{f=\frac{1}{T}} and radians/sec for \inlineTFinmath{\omega
=\frac{2\pi }{T}}.



\Exercise{Problem 1}
\ExerciseText{Obtain the frequency of the following sinusoids:

sin( 4\(\pi \)t), y sin(6 t).}
\begin{CellGroup}
\Subsection*{Real and complex exponential signals}
Another basic signal is the continuous-time exponential. Later we will see that real exponentials arise naturally in the analysis of electric circuits,
while complex exponentials greatly simplify fundamental mathematical formulations of linear signal theory. A  complex exponential is actually a complex
valued function, whose real and imaginary parts are sinusoidal. Recall from Euler's identity that

\inlineTFinmath{\multsp {{\ExponentialE }^{\multsp j\multsp {{\omega }_0}t}}\multsp =\multsp \cos({{\omega }_0}t)+\multsp j\multsp \sin({{\omega
}_0}t)}

 The fundamental period \inlineTFinmath{T}, of the signal \inlineTFinmath{x(t)={{\ExponentialE }^{\multsp j\multsp {{\omega }_0}t}}} is therefore
given by (2). Converting between complex exponentials and trigonometric functions is simple with {\itshape Mathematica}.



\begin{CellGroup}
\dispSFinmath{\Mfunction{ExpToTrig}[\exp [\imag \multsp {{\omega }_0}t]]}
\dispSFoutmath{\cos [t\multsp {{\omega }_0}]+\ImaginaryI \multsp \sin [t\multsp {{\omega }_0}]}

\end{CellGroup}
\begin{CellGroup}
\dispSFinmath{\Mfunction{TrigToExp}[\%]}
\dispSFoutmath{{{\ExponentialE }^{\ImaginaryI \multsp t\multsp {{\omega }_0}}}}

\end{CellGroup}
Here we plot the real and complex exponentials. A real valued exponential is a rapidly decreasing (increasing) and unbounded function. The complex
exponential is a periodic complex function. There are many equivalent ways of plotting complex functions, here we show a combined plot of the real
(in red) and imaginary (in blue) parts of the function.

\begin{CellGroup}
\dispSFinmath{\MathBegin{MathArray}{l}
\Mvariable{DisplayTogetherArray}[  \\
\noalign{\vspace{0.5ex}}
\hspace{1.em} \{\Mfunction{Plot}[\{\Mfunction{Re}[\exp [-\multsp \imag \multsp t]],\Mfunction{Im}[\exp [-\multsp \imag \multsp t]]\},\{t,0,2\multsp
\pi \},\Mvariable{PlotLabel}\rightarrow   \\
\noalign{\vspace{0.5ex}}
\hspace{4.em} \Mvariable{Se\NTilde al\multsp Compleja},\Mvariable{PlotStyle}->\{\Mfunction{RGBColor}[1,0,0],\Mfunction{RGBColor}[0,0,1]\}],  \\
\noalign{\vspace{0.5ex}}
\hspace{2.em} \Mfunction{Plot}[\exp [-\multsp t],\{t,0,2\multsp \pi \},\Mvariable{PlotLabel}\rightarrow \Mvariable{Se\NTilde al\multsp Real}]\}];\\
\MathEnd{MathArray}}
\mathGraphic{lecture02.tex_gr2.eps}

\end{CellGroup}
\begin{CellGroup}
\Exercise{Problem 2}
\ExerciseText{Use Euler's identity to convert the following sum of complex exponentials to trigonometric form:

    \inlineTFinmath{{{\ExponentialE }^{\multsp j\multsp {{\omega }_0}t}}\multsp +\multsp \multsp {{\ExponentialE }^{-j\multsp {{\omega }_0}t\multsp
}}}

 Do this problem by hand and with {\itshape Mathematica}.}



\end{CellGroup}

\end{CellGroup}
\begin{CellGroup}
\Subsection*{Sinc function}
The final signal presented in this section is the so-called sinc function, \inlineTFinmath{\Muserfunction{sinc}(t)}, defined for real \inlineTFinmath{t}
as



\dispSFinmath{\Muserfunction{sinc}[\Mvariable{t\_}]:=\frac{\sin [t]}{t}}
The sinc function is of fundamental importance in sampling theory and in the theory of generalized functions. It also shows up as the result of calculating
the inverse, continuous-time Fourier transform of an ideal lowpass filter.

\begin{CellGroup}
\dispSFinmath{\MathBegin{MathArray}{l}
\Mfunction{Plot}[\Muserfunction{sinc}[\pi \multsp t],\{t,-2\multsp \pi ,2\multsp \pi \},  \\
\noalign{\vspace{0.5ex}}
\hspace{1.em} \Mvariable{PlotRange}->\Mvariable{All},\Mvariable{PlotLabel}\rightarrow \Mvariable{Funcion\multsp sinc(t)}];\\
\MathEnd{MathArray}}
\mathGraphic{lecture02.tex_gr3.eps}

\end{CellGroup}
\begin{CellGroup}
\Exercise{Problem 3}
\ExerciseText{Find the zeros of \inlineTFinmath{\Muserfunction{sinc}(\pi \multsp t)}.}



\end{CellGroup}

\end{CellGroup}

\end{CellGroup}
\begin{CellGroup}
\Section*{2 Discrete Signals}


Signals may also be sequences. Sequences, in contrast to continuous-time  signals, are defined over an infinite, semi-infinite or finite set of integers.
A discrete-time signal is therefore represented by a list of real or complex values.

\begin{CellGroup}
\Subsection*{Real-valued sinusoidal sequences}
Following the examples in Section 1, we begin by plotting two discrete-time sinusoids. Do these sequences look sinusoidal? Are they indeed periodic?
These are some of the questions we will address in this section.

Here we define the two sinusoidal sequences.

\dispSFinmath{\MathBegin{MathArray}{l}
\Mvariable{sig1}=\Mfunction{Table}[\sin [0.589\multsp n],\{n,0,31\}];  \\
\noalign{\vspace{0.5ex}}  \\
\Mvariable{sig2}=\Mfunction{Table}[\sin [1.767\multsp n],\{n,0,31\}];
\MathEnd{MathArray}}
Here we plot the resulting sequences.

\begin{CellGroup}
\dispSFinmath{\Muserfunction{DisplayTogetherArray}[\{\Muserfunction{LinePlot}[\Mvariable{sig1}],\Muserfunction{LinePlot}[\Mvariable{sig2}]\}];}
\mathGraphic{lecture02.tex_gr4.eps}

\end{CellGroup}

\end{CellGroup}
\begin{CellGroup}
\Subsection*{Real and complex exponential sequences}
The general discrete-time exponential signal is of the form 

	\inlineTFinmath{x[n]\multsp =\multsp {r^n}}, 



where the value \inlineTFinmath{r} may be real or complex. For real values of {\itshape r} the signal \inlineTFinmath{x[n]} is known as the geometric
sequence, while for complex values of the form, \inlineTFinmath{r=\multsp {{\ExponentialE }^{\multsp j\multsp \omega \multsp n\multsp }}}, \inlineTFinmath{x[n]}
is known as the discrete-time complex exponential. Here is a plot of a geometric sequence.



\begin{CellGroup}
\dispSFinmath{\Muserfunction{LinePlot}[\Mfunction{Table}[{{0.77}^n},\{n,0,32\}],\Mvariable{PlotRange}->\Mvariable{All},\Mvariable{PlotLabel}\rightarrow
\Mstring{${{0.77}^n}$}
];}
\mathGraphic{lecture02.tex_gr5.eps}

\end{CellGroup}
Later in the semester we shall make considerable use of the complex valued  discrete-time exponential sequence. Here, we first define a complex exponential
signal.

\dispSFinmath{x[\Mvariable{n\_}]:=\exp [\imag \multsp 0.589\multsp n]}
 Here are plots of the real and imaginary parts of \inlineTFinmath{x[n]}. 



\begin{CellGroup}
\dispSFinmath{\MathBegin{MathArray}{l}
\Mvariable{DisplayTogetherArray}[  \\
\noalign{\vspace{0.5ex}}
\hspace{1.em} \{\Muserfunction{LinePlot}[\Mfunction{Table}[\Mfunction{Re}[x[n]],\{n,0,32\}],\Mvariable{PlotLabel}\rightarrow \Mvariable{Parte\multsp
real}],  \\
\noalign{\vspace{0.5ex}}
\hspace{2.em} \Muserfunction{LinePlot}[\Mfunction{Table}[\Mfunction{Im}[x[n]],\{n,0,32\}],\Mvariable{PlotLabel}\rightarrow \Mvariable{Parte\multsp
imaginaria}]\}];\\
\MathEnd{MathArray}}
\mathGraphic{lecture02.tex_gr6.eps}

\end{CellGroup}

\end{CellGroup}
\begin{CellGroup}
\Subsection*{Periodicity Properties of Discrete Time Complex Exponentials}
Discrete-time exponential of the form \inlineTFinmath{{{\ExponentialE }^{\multsp j\multsp \omega \multsp n\multsp }}} is periodic with period 2\(\pi
\), namely



\inlineTFinmath{{{\ExponentialE }^{\multsp \multsp j\multsp (\omega \multsp \multsp +\multsp 2\pi )\multsp n}}\multsp =\multsp \multsp {{\ExponentialE
}^{\multsp \multsp j\multsp \omega \multsp n}}}.

\begin{CellGroup}
\Exercise{Problem 4}
\ExerciseText{Prove (4).}

\end{CellGroup}
Consequently, the signal \inlineTFinmath{x[n]\multsp =\multsp {{\ExponentialE }^{\multsp j\multsp {{\omega }_0}\multsp n}}} does NOT have a continually
increasing rate of oscillation as the  frequency \inlineTFinmath{{{\omega }_0}}is increased (verify this by plotting \inlineTFinmath{x[n]} for a
few different frequencies in the range \inlineTFinmath{\omega \leq {{\omega }_0}\leq \omega +2\pi }). Note that this property makes the discrete-time
complex exponential significantly different from the continuous-time complex exponential (again, you can easily verify this by plotting \inlineTFinmath{x(t)}
and \inlineTFinmath{x[n]}, an comparing for the same fundamental frequency values).

The discrete time exponential of the form \inlineTFinmath{{{\ExponentialE }^{\multsp j\multsp {{\omega }_0}\multsp n\multsp }}} is also periodic
with period \inlineTFinmath{N}, namely,



\inlineTFinmath{{{\ExponentialE }^{\multsp \multsp j\multsp {{\omega }_0}\multsp (n+\multsp N)}}\multsp =\multsp \multsp {{\ExponentialE }^{\multsp
\multsp j\multsp {{\omega }_0}n}}},

if-and-only-if 

\inlineTFinmath{\frac{{{\omega }_0}}{2\pi }\multsp =\multsp \multsp \frac{m}{N}}, 	where \inlineTFinmath{m,\multsp N\multsp \epsilon \multsp \DoubleStruckCapitalZ
}

For (6) to be satisfied, the ratio \inlineTFinmath{\frac{{{\omega }_0}}{2\pi }} must be a rational number.



\begin{CellGroup}
\Exercise{Problem 5}
\ExerciseText{Determine if the following signals are periodic with respect to {\itshape n}:

\inlineTFinmath{{x_1}[n]=\multsp {{\ExponentialE }^{\multsp j\multsp 4\pi \multsp n}}},	\inlineTFinmath{{x_2}[n]=\multsp {{\ExponentialE }^{\multsp
j\multsp 6\multsp n}}}}



\end{CellGroup}

\end{CellGroup}

\end{CellGroup}
\begin{CellGroup}
\Section*{3 Special Signals}


There is a general class of functions called generalized functions. These are not ordinary function and are of significant importance in the theoretical
 analysis of linear systems. Two members of this class of functions are  especially important: the ideal impulse (known as the Dirac delta) and the
 step function. The Dirac delta function is an important tool in the analysis  of the input-output behavior of linear, time-invariant systems. It
also  plays an important role in Fourier theory.

\begin{CellGroup}
\Subsection*{Unit Impulse Signal}
The Dirac delta function \inlineTFinmath{\delta (t)} is the {\itshape singular} function such that:



\inlineTFinmath{\int _{-\infty }^{\infty }\delta (t)\multsp f(t)\DifferentialD t\multsp =\multsp f(0)}

Singular functions are not defined "pointwise" as ordinary continuous functions, but only "under" an integral sine. It is their effect on other functions
that defines them. (7) is also known as the {\itshape sampling} property of the Dirac delta.



Here we demonstrate the sampling property of the impulse function.

\dispSFinmath{\Mfunction{If}[\$VersionNumber\neq 4.,<<\Mvariable{Calculus`DiracDelta`}]}
\begin{CellGroup}
\dispSFinmath{\int _{-\infty }^{\infty }\Mfunction{DiracDelta}[t]\multsp f[t]\multsp \DifferentialD t}
\dispSFoutmath{f[0]}

\end{CellGroup}
Additional properties of \inlineTFinmath{\delta (t)}:



\inlineTFinmath{a\multsp \delta (t)\multsp +\multsp b\multsp \delta (t)\multsp =\multsp (a+b)\multsp \delta (t)}		(* linearity *)

\inlineTFinmath{\int _{-\infty }^{\infty }\delta (t-{t_0})\multsp f(t)\DifferentialD t\multsp =\multsp f({t_0})}		(* time translation *)

\inlineTFinmath{\int _{-\infty }^{\infty }\delta (t)\multsp f(t)\DifferentialD t\multsp =\multsp \int _{-\infty }^{\infty }\multsp \delta (t)\multsp
f(0)\DifferentialD t}		(* multiplication by a function *)

\begin{CellGroup}
\Exercise{Problem 6}
\ExerciseText{Obtain the value of the following integral: \inlineTFinmath{\int _{-\infty }^{\infty }\delta (t)\multsp \DifferentialD t\multsp }.
Prove your result and verify using {\itshape Mathematica}. }



\end{CellGroup}

\end{CellGroup}
Here are two common approximations of the Dirac delta.

\begin{CellGroup}
\dispSFinmath{\MathBegin{MathArray}{l}
\Mfunction{Plot}\big[\frac{1}{{\sqrt{\epsilon \multsp \pi }}}\exp \big[-\frac{{t^2}}{\epsilon }\big]\multsp /.\multsp \epsilon ->0.001,\multsp \{t,-1,1\},\multsp
  \\
\noalign{\vspace{1.16667ex}}
\hspace{1.em} \Mvariable{PlotRange}->\Mvariable{All},\Mvariable{PlotLabel}\rightarrow \Mvariable{Aproximacion\multsp del\multsp impulso\multsp unitario}\big];\\
\MathEnd{MathArray}}
\mathGraphic{lecture02.tex_gr7.eps}

\end{CellGroup}
\begin{CellGroup}
\dispSFinmath{\MathBegin{MathArray}{l}
\Mfunction{Plot}\big[\frac{\sin [w\multsp t]}{\pi \multsp t}\multsp /.\multsp w->100,\multsp \{t,-1,1\},\Mvariable{PlotRange}->\Mvariable{All}, 
\\
\noalign{\vspace{1.07292ex}}
\hspace{1.em} \Mvariable{PlotLabel}\rightarrow \Mvariable{Aproximacion\multsp del\multsp impulso\multsp unitario}\big];\\
\MathEnd{MathArray}}
\mathGraphic{lecture02.tex_gr8.eps}

\end{CellGroup}
\begin{CellGroup}
\Subsection*{Unit Step Signal}
 Here is a second important generalized function, the so-called unit step function, \inlineTFinmath{u(t)}. The unit step function is defined:



\inlineTFinmath{u(t)\multsp =\multsp \big\{\MathBegin{MathArray}[c]{c}
	0,\multsp \multsp \multsp \multsp \multsp \multsp t<0 \\
	1,\multsp \multsp \multsp \multsp \multsp \multsp \multsp t\geq 0 
  \MathEnd{MathArray}}

\begin{CellGroup}
\dispSFinmath{\MathBegin{MathArray}{l}
\Mfunction{Plot}[\Mfunction{UnitStep}[t],\{t,-2,5\},\Mvariable{PlotRange}->\{0,1.5\},  \\
\noalign{\vspace{0.5ex}}
\hspace{1.em} \Mvariable{PlotLabel}->\Mvariable{Escalon\multsp unitario},\multsp \Mvariable{PlotStyle}->\Mfunction{RGBColor}[1,0,0]];\\
\MathEnd{MathArray}}
\mathGraphic{lecture02.tex_gr9.eps}

\end{CellGroup}
The unit step and impulse are closely related. In fact, the unit step is a running integral of the unit impulse. 

\inlineTFinmath{u(t)\multsp =\multsp \int _{-\infty }^{t}\delta (\tau )\DifferentialD \tau }

Conversely, the unit  impulse is a derivative of the unit step. {\itshape Mathematica} "understands" this property, as can be easily verified.



\begin{CellGroup}
\dispSFinmath{{{\partial }_t}\multsp \Mfunction{UnitStep}[t]}
\dispSFoutmath{\Mfunction{DiracDelta}[t]}

\end{CellGroup}
\Subsection*{Unit Sample and Unit Step Sequences}
The discrete-time equivalents of these two special signals are the so-called  Kronecker delta, \inlineTFinmath{\delta [n]} (also known as the unit
sample) and the unit step sequence, \inlineTFinmath{u[n]}. These are however ordinary signals. The two definitions are:



\inlineTFinmath{\delta [n]\multsp =\multsp \big\{\MathBegin{MathArray}[c]{c}
	1,\multsp \multsp \multsp n=0 \\
	0,\multsp n\neq 0 
  \MathEnd{MathArray}}

and

\inlineTFinmath{u[n]\multsp =\multsp \big\{\MathBegin{MathArray}[c]{c}
	0,\multsp \multsp \multsp n<0 \\
	1,\multsp n\geq 0 
  \MathEnd{MathArray}}

Similar relationships exist between the unit sample and unit step sequences as for their continuous-time counterparts. The unit sample is a {\itshape
first-difference} of the discrete-time step:



\inlineTFinmath{\delta [n]=\multsp u[n]-u[n-1]}

and the unit step is a {\itshape running sum} of the unit sample.



\inlineTFinmath{u[n]\multsp =\multsp \sum _{m=-\infty }^{n}\delta [m]}

Here we plot the unit sample (shifted) and step sequences.

\begin{CellGroup}
\dispSFinmath{\MathBegin{MathArray}{l}
\Mvariable{DisplayTogetherArray}[  \\
\noalign{\vspace{0.5ex}}
\hspace{1.em} \{\Muserfunction{LinePlot}[\Mfunction{Table}[\{n,\Mfunction{KroneckerDelta}[n-3]\},\{n,-3,12\}],  \\
\noalign{\vspace{0.5ex}}
\hspace{3.em} \Mvariable{PlotLabel}->\Mvariable{impulso\multsp unitario},\Mvariable{PointStyle}\rightarrow \Mfunction{PointSize}[0.03]],  \\
\noalign{\vspace{0.5ex}}\hspace{1em}\hspace{1em}\hspace{1em}\Muserfunction{LinePlot}[\Mfunction{Table}[\{n,\Mfunction{UnitStep}[n]\},\{n,-3,12\}],
 \\
\noalign{\vspace{0.5ex}}
\hspace{3.em} \Mvariable{PlotLabel}->\Mvariable{escalon\multsp unitario},\Mvariable{PointStyle}\rightarrow \Mfunction{PointSize}[0.03]]\}];\\
\MathEnd{MathArray}}
\mathGraphic{lecture02.tex_gr10.eps}

\end{CellGroup}
Note that any sequence can be represented as a weighted sum of unit samples. For example consider the sequence \inlineTFinmath{x[n]} defined below,




\inlineTFinmath{x[n]\multsp =\multsp \big\{\MathBegin{MathArray}[c]{l}
	1,\multsp n=0 \\
	2,\multsp n=1,2 \\
	-1,\multsp n=3 \\
	0,\multsp \Mvariable{otherwise} 
  \MathEnd{MathArray}\multsp }

This sequence may be written in the form

\inlineTFinmath{x[n]=\multsp \delta [n]\multsp +\multsp 2\multsp \delta [n-1]\multsp +\multsp 2\multsp \delta [n-2]\multsp \multsp -\multsp \delta
[n-3]}


\end{CellGroup}

\end{CellGroup}
\begin{CellGroup}
\Section*{PROBLEM SOLUTIONS}
\begin{CellGroup}
\Exercise{Problem 1}
\ExerciseText{Obtain the frequency of the following sinusoids:

    \inlineTFinmath{\sin(\multsp 4\multsp \pi \multsp t)},	\inlineTFinmath{\cos(\multsp 6\multsp t)}.}


\ExerciseText{{\bfseries Solution}.

 \inlineTFinmath{\sin(4\pi \multsp t)\multsp =\multsp \sin(4\pi \multsp t+2\pi )\multsp =\multsp \sin\Big(4\pi \Big(\multsp t+\frac{1}{2}\Big)\Big)},
	thus \inlineTFinmath{T=\frac{1}{2}}

 \inlineTFinmath{\cos(6\multsp t)\multsp =\multsp \cos(6\multsp t+2\pi )\multsp =\multsp \cos\big(6\big(t+\frac{\pi }{3}\big)\big)},		thus \inlineTFinmath{T=\frac{\pi
}{3}} }



\end{CellGroup}
\begin{CellGroup}
\Exercise{Problem 2}
\ExerciseText{Use Euler's identity to convert the following sum of complex exponentials  to trigonometric form:

    \inlineTFinmath{{{\ExponentialE }^{\multsp j\multsp {{\omega }_0}t}}\multsp +\multsp \multsp {{\ExponentialE }^{-j\multsp {{\omega }_0}t}}}

 Do this problem both by hand and with {\itshape Mathematica}.}


\ExerciseText{{\bfseries Solution.} Using well know rules of exponents, we get  

\inlineTFinmath{\MathBegin{MathArray}{l}
{{\ExponentialE }^{\multsp \ImaginaryI \multsp {{\omega }_0}\multsp \multsp t}}\multsp +\multsp {{\ExponentialE }^{\multsp -\ImaginaryI \multsp {{\omega
}_0}\multsp \multsp t}}=\multsp   \\
\noalign{\vspace{0.666667ex}}
\hspace{1.em} \cos({{\omega }_0}t)+\multsp \ImaginaryI \multsp \sin({{\omega }_o}t)+\cos({{\omega }_0}t)-\multsp \ImaginaryI \multsp \sin({{\omega
}_o}t)=  \\
\noalign{\vspace{0.666667ex}}
\hspace{2.em} 2\multsp \cos({{\omega }_0}t)\\
\MathEnd{MathArray}}}


\begin{CellGroup}
\dispSFinmath{\exp [\imag \multsp {{\omega }_0}\multsp t]+\exp [-\imag \multsp {{\omega }_0}\multsp t]//\Mfunction{ExpToTrig}\multsp //\Mfunction{Simplify}}
\dispSFoutmath{2\multsp \cos [t\multsp {{\omega }_0}]}

\end{CellGroup}

\end{CellGroup}
\begin{CellGroup}
\Exercise{Problem 3}
\ExerciseText{Find the zeros of \inlineTFinmath{\Muserfunction{sinc}(\pi \multsp t)}.}


\ExerciseText{{\bfseries Solution.} The zeros are the same as zeros of \inlineTFinmath{\sin(\pi \multsp t)}, so

\inlineTFinmath{\pi \multsp t\multsp =\multsp \pi \multsp k,} 	and therefore

\inlineTFinmath{\multsp t=\multsp k},		for	\inlineTFinmath{\multsp k\multsp =\multsp 0,\multsp \pm 1,\multsp \pm 2,\multsp ...}}



\end{CellGroup}
\begin{CellGroup}
\Exercise{Problem 4}
\ExerciseText{Prove (4).}
\ExerciseText{{\bfseries Solution.}}


\ExerciseText{\inlineTFinmath{{{\ExponentialE }^{\multsp \multsp j\multsp (\omega \multsp \multsp +\multsp 2\pi )\multsp n}}\multsp =\multsp \multsp
{{\ExponentialE }^{\multsp \multsp j\multsp \omega \multsp n}}\multsp {{\ExponentialE }^{\multsp \multsp j\multsp 2\pi \multsp n}}\multsp =\multsp
\multsp {{\ExponentialE }^{\multsp \multsp j\multsp \omega \multsp n}}\multsp *\multsp 1=\multsp {{\ExponentialE }^{\multsp \multsp j\multsp \omega
\multsp n}}}}



\end{CellGroup}
\begin{CellGroup}
\Exercise{Problem 5}
\ExerciseText{Determine if the following signals are periodic:

\inlineTFinmath{{x_1}[n]=\multsp {{\ExponentialE }^{\multsp j\multsp 4\pi \multsp n}}},	\inlineTFinmath{{x_2}[n]=\multsp {{\ExponentialE }^{\multsp
j\multsp 6\multsp n}}}}


\ExerciseText{{\bfseries Solution.}}


\ExerciseText{\inlineTFinmath{{x_1}[n]}:	\inlineTFinmath{\frac{{{\omega }_0}}{2\pi }=\frac{4\pi }{2\pi }=\multsp 2},	therefore \inlineTFinmath{{x_1}[n]}
is periodic

\inlineTFinmath{{x_2}[n]}: 	\inlineTFinmath{\frac{{{\omega }_0}}{2\pi }=\frac{6}{2\pi }}, 		therefore \inlineTFinmath{{x_2}[n]} is not periodic}



\end{CellGroup}
\begin{CellGroup}
\Exercise{Problem 6}
\ExerciseText{Obtain the value of the following integral: \inlineTFinmath{\int _{-\infty }^{\infty }\delta (t)\multsp \DifferentialD t\multsp }.
Prove your result and verify using {\itshape Mathematica}. }


\ExerciseText{{\bfseries Solution}. We get the result immediately from (7) by taking \inlineTFinmath{f(t)=1}. Also,}


\begin{CellGroup}
\dispSFinmath{\int _{-\infty }^{\infty }\Mfunction{DiracDelta}[t]\DifferentialD t}
\dispSFoutmath{1}

\end{CellGroup}

\end{CellGroup}

\end{CellGroup}
\Section*{HOMEWORK PROBLEMS}

\end{document}
